\clearpage\begingroup
\SourceLeaf{131}\NoSourceHead
\begin{center}
[III.]\par\medskip
\Large DIOPTRIK.\par\medskip
\rule{.25\linewidth}{.4pt}\par\medskip
{\normalsize [Aus Scheda An, Cereri, Palladi, Iunoni sacrum, Febr. 1805, S. 70, 71.}\par\medskip
\rule{.25\linewidth}{.4pt}
\end{center}
\noindent\hspace*{\parindent}Drei Gläser, deren Brennweiten [in Fuss]
\[
2{,}090\left[=\frac1f\right],\quad 0{,}353\left[=\frac1{f'}\right],\quad 0{,}152\left[=\frac1{f''}\right]
\]
\noindent [sind,] sollen für ein kurzsichtiges Auge, dessen Gesichtsweite $=0{,}4$ ist, zu einem Fernrohre arrangirt werden in welchem der farbige Rand gehoben ist.
\begin{center}[I\textsuperscript{te} Hypothese.]\end{center}
\[
h=1{,}937,\qquad h'=0{,}2079
\]
\[
\begin{aligned}
\alpha&=1\\
\beta&=0{,}47847\\
\alpha'&=-0{,}07321\\
\beta'&=-0{,}68584\\
\alpha''&=-0{,}06938\\
\beta''&=+0{,}22939
\end{aligned}\qquad\qquad
\begin{aligned}
\gamma&=0\\
\delta&=1\\
\gamma'&=1{,}937\\
\delta'&=4{,}4867\\
\gamma''&=-1{,}0042\\
\delta''&=-11{,}0933.
\end{aligned}
\]
\begin{center}[Berechnung des farbigen Randes:]\end{center}
\[
\partial\delta'\left[=\gamma'f'\frac{\partial n}{n-1}\right]
=5{,}4867\mathbin{.}\left[\frac{\partial n}{n-1}\right],
\]
\[
\partial\gamma''\left[=h'\gamma'f'\frac{\partial n}{n-1}
=h'(\delta+\delta')\frac{\partial n}{n-1}\right]
=1{,}1407\mathbin{.}\left[\frac{\partial n}{n-1}\right],
\]
\[
\partial\delta''\left[=\{(f''h'-1)(\delta'+\delta)+(\delta''+\delta')\}\frac{\partial n}{n-1}\right]
=-4{,}5887\mathbin{.}\left[\frac{\partial n}{n-1}\right],
\]
\clearpage\SourceLeaf{132}\HistoricalHead{127}{R}{DIOPTRIK.}
\[
\begin{array}{r@{\quad}l}
[\alpha''\partial\delta''={}]&+0{,}31836\mathbin{.}\left[\frac{\partial n}{n-1}\right]\\[5pt]
{}[-\beta''\partial\gamma''={}]&\phantom{+}0{,}26166\mathbin{.}\left[\frac{\partial n}{n-1}\right]\\[5pt]\hline
{}[\alpha''\partial\delta''-\beta''\partial\gamma''={}]&+0{,}0567\mathbin{.}\left[\frac{\partial n}{n-1}\right].
\end{array}
\]
\begin{center}[II\textsuperscript{te} Hypothese.]\end{center}
\[
h=1{,}947,\qquad h'=0{,}2036
\]
\[
\begin{aligned}
\alpha&=1\\
\beta&=0{,}47847\\
\alpha'&=-0{,}06842\\
\beta'&=-0{,}67230\\
\alpha''&=-0{,}06846\\
\beta''&=+0{,}22190
\end{aligned}\qquad\qquad
\begin{aligned}
\gamma&=0\\
\delta&=1\\
\gamma'&=1{,}947\\
\delta'&=4{,}5156\\
\gamma''&=-1{,}0276\\
\delta''&=-11{,}2781.
\end{aligned}
\]
\begin{center}[Berechnung des farbigen Randes:]\end{center}
\[
\begin{aligned}
\partial\delta'&=5{,}5156\mathbin{.}\left[\frac{\partial n}{n-1}\right],\\[4pt]
\partial\gamma''&=1{,}1230\mathbin{.}\left[\frac{\partial n}{n-1}\right],\\[4pt]
\partial\delta''&=-4{,}8899\mathbin{.}\left[\frac{\partial n}{n-1}\right],\\[4pt]
{}[\alpha''\partial\delta''={}]&\ +0{,}33476\mathbin{.}\left[\frac{\partial n}{n-1}\right],\\[4pt]
{}[-\beta''\partial\gamma''={}]&\ -0{,}24919\mathbin{.}\left[\frac{\partial n}{n-1}\right],\\[4pt]
{}[\alpha''\partial\delta''-\beta''\partial\gamma''={}]&\ +0{,}0856\mathbin{.}\left[\frac{\partial n}{n-1}\right].
\end{aligned}
\]
\begin{center}III\textsuperscript{te} Hypothese.\end{center}
\[
[h={}]\ 1{,}917,\qquad [h'={}]\ 0{,}2165
\]
\[
\begin{aligned}
[\alpha&=1\\
\beta&=0{,}47847\\
\alpha'&=-0{,}08277\\
\beta'&=-0{,}71295\\
\alpha''&=-0{,}07158\\
\beta''&=+0{,}24203
\end{aligned}\qquad\qquad
\begin{aligned}
\gamma&=0\\
\delta&=1\\
\gamma'&=1{,}917\\
\delta'&=4{,}43059\\
\gamma''&=-0{,}95778\\
\delta''&=-10{,}73177.
\end{aligned}
\]
\clearpage\SourceLeaf{133}\HistoricalHead{128}{L}{PHYSIK.\quad NACHLASS.}
\begin{center}Berechnung des farbigen Randes:\end{center}
\[
\begin{aligned}
\partial\delta'&=5{,}4306\mathbin{.}\frac{\partial n}{n-1},\\[4pt]
\partial\gamma''&=1{,}1757\mathbin{.}\frac{\partial n}{n-1},\\[4pt]
\partial\delta''&=-3{,}99698\mathbin{.}\frac{\partial n}{n-1},\\[4pt]
\alpha''\partial\delta''&=0{,}2861\mathbin{.}\frac{\partial n}{n-1},\\[4pt]
-\beta''\partial\gamma''&=-0{,}2883\mathbin{.}\frac{\partial n}{n-1},\\[4pt]
\text{farbiger Rand}\quad\alpha''\partial\delta''-\beta''\partial\gamma''&=-0{,}0024\mathbin{.}\frac{\partial n}{n-1}].
\end{aligned}
\]
\noindent\hspace*{\parindent}Ort des Auges $0{,}089$\par
\noindent\hspace*{\parindent}Blendung v[om] Objectiv $0{,}116$\par
\noindent\hspace*{\parindent}Gesichtsfeld $4^{0}\,18'$.
\begin{center}\rule{.25\linewidth}{.4pt}\end{center}
\clearpage\endgroup
